\chapter{Conclusiones y trabajo futuro}
\label{cap:conclusiones}

\section{Conclusiones}

\instruccion{Redacte un cierre integrado y sustentado. Responda de forma directa
la pregunta principal y las preguntas específicas; establezca la conclusión
respecto de la hipótesis; explique en qué medida se alcanzaron los objetivos;
identifique las contribuciones realmente obtenidas; delimite qué puede afirmarse
con la evidencia disponible; reconozca las limitaciones del estudio; e indique,
cuando corresponda, la disponibilidad, versión, licencia o restricciones de
datos, código, modelos, instrumentos y otros artefactos. No introduzca resultados
nuevos ni repita el resumen de cada capítulo.}

Sustituya este párrafo por las conclusiones de la tesis.

\section{Trabajo futuro}

\instruccion{Proponga líneas concretas derivadas de los hallazgos y las
limitaciones. Para cada una, explique qué pregunta, evidencia, comparación o
capacidad ampliaría. Evite presentar como investigación futura una lista de
funcionalidades pendientes sin relación con el conocimiento generado.}

Sustituya este párrafo por las líneas de trabajo futuro.

